
\documentclass{article} % For LaTeX2e
\usepackage[submission]{colm2026_conference}

% Only load packages not already provided by colm2026_conference.sty
\usepackage{booktabs}
\usepackage{amsmath}
\usepackage{amssymb}
\usepackage{graphicx}
\usepackage{algorithm}
\usepackage{algorithmic}
\usepackage{multirow}
\usepackage[table]{xcolor}

\definecolor{darkblue}{rgb}{0, 0, 0.5}
\definecolor{darkgreen}{rgb}{0.0, 0.5, 0.0}

\newcommand{\CATPO}{\textsc{CATPO}}
\newcommand{\TreeRPO}{\textsc{TreeRPO}}
\newcommand{\placeholder}[1]{{\color{red}\textbf{[#1]}}}

\title{CATPO: Critique-Augmented Tree Policy Optimization}

\author{Author Name \\
Department \\
Institution \\
\texttt{email@institution.edu}
}

\begin{document}

\maketitle

% ============================================================================
% ABSTRACT
% ============================================================================
\begin{abstract}
Reinforcement learning with verifiable rewards (RLVR) has become a dominant paradigm for training large language models (LLMs) to reason. Recent tree-based policy optimization methods such as \TreeRPO extend flat trajectory sampling with tree-structured rollouts, enabling the LLM policy to receive dense, step-level advantage estimates without training a separate process reward model. However, not all trees are equally informative for updating the policy: trees where all leaves succeed (\emph{dead-correct}), all leaves fail (\emph{dead-wrong}), or the policy already predicts the reward distribution (\emph{stale}) contribute little to the policy gradient, wasting compute. We introduce \textbf{\CATPO} (Critique-Augmented Tree Policy Optimization), which addresses this through two complementary mechanisms: (1)~a \emph{tree fitness function} $F(T) = \mathcal{H}(\mathbf{r}) \cdot (1 - \rho_{\pi,r})$ that quantifies each tree's informativeness for policy improvement at zero extra compute, combining the entropy of leaf outcomes with the policy's surprise at the reward distribution; and (2)~a \emph{critique-guided healing} procedure that identifies the shallowest failing point in dead-wrong trees, generates natural-language critiques, and grafts refined continuations to recover training signal for the policy. \CATPO weights the policy gradient loss by tree fitness, so that informative trees dominate parameter updates while healed trees naturally contribute through the standard \TreeRPO objective after re-propagation. Experiments on Qwen2.5-Math-1.5B trained with the MATH dataset show that \CATPO achieves \placeholder{XX.X}\% macro accuracy across four benchmarks (AIME24, MATH-500, OlympiadBench, Minerva), improving over \TreeRPO by \placeholder{+X.X}\% and GRPO by \placeholder{+X.X}\%, while reducing average response length by \placeholder{XX.X}\%.
\end{abstract}


% ============================================================================
% 1. INTRODUCTION
% ============================================================================
\section{Introduction}
\label{sec:intro}

Recent breakthroughs in reinforcement learning have brought transformative progress to large language model (LLM) reasoning. Models such as OpenAI O1~\citep{o1}, DeepSeek-R1~\citep{guo2025deepseek}, and QwQ~\citep{qwq} achieve remarkable performance on complex reasoning tasks by training the LLM policy with Reinforcement Learning with Verifiable Rewards (RLVR), where the model continuously explores reasoning paths toward correct answers on problems with verifiable solutions.

A central challenge in training the LLM policy via RLVR is \emph{credit assignment}: trajectory-level rewards (whether from outcome reward models or rule-based verification) provide limited guidance for determining which intermediate reasoning steps actually contributed to the final outcome. This has motivated two lines of work to obtain finer-grained advantage estimates for the policy gradient. The first trains explicit process reward models (PRMs)~\citep{verify_stepbystep, mathshepherd, omegaprm} --- separate networks that score each reasoning step --- which require expensive step-level annotation. The second avoids training any auxiliary model and instead uses tree-structured sampling: by expanding multiple continuations at each reasoning step, methods like \TreeRPO~\citep{treerpo} and TreeRL~\citep{treerl} estimate step-level advantages directly from leaf outcomes propagated bottom-up through the tree, providing dense policy gradient signal without a separate reward model.

However, tree-based methods treat \emph{all trees equally} in the training loss. In practice, a large fraction of sampled trees are uninformative:
\begin{itemize}
    \item \textbf{Dead-correct trees} where all leaves produce the right answer --- the model already solves this problem, so no contrastive signal exists.
    \item \textbf{Dead-wrong trees} where all leaves fail --- no branch provides a positive example, so the group-relative advantage yields zero effective policy gradient.
    \item \textbf{Stale trees} where the policy's confidence already aligns with the reward distribution --- the model has internalized this pattern, and further training on it yields diminishing returns.
\end{itemize}

This waste is not merely theoretical. DAPO~\citep{dapo} introduced Dynamic Sampling to filter zero-variance \emph{prompts} (all rollouts correct or all incorrect), and GRESO~\citep{greso} predicts uninformative prompts before rolling out. But these operate at the prompt level, not the tree level: the same prompt can produce an informative tree or a dead one depending on the sampling, and the internal structure of the tree (which branches the policy finds surprising) carries signal that prompt-level filtering discards.

We propose \textbf{\CATPO} (Critique-Augmented Tree Policy Optimization), which introduces two mechanisms to address tree-level sample inefficiency:

\paragraph{1. Tree Fitness Scoring.} We define a fitness function $F(T) = \mathcal{H}(\mathbf{r}) \cdot (1 - \rho_{\pi,r})$ that scores each tree's informativeness by combining the binary entropy of leaf outcomes (reward diversity) with the Pearson decorrelation between the policy's log-probabilities and propagated rewards (policy surprise). $F(T)$ is computed in closed form from quantities already available after tree sampling, requiring zero additional forward passes. We use $F(T)$ to weight the training loss, concentrating gradient updates on informative trees.

\paragraph{2. Critique-Guided Healing.} For dead-wrong trees --- where all branches fail and $F \approx 0$ --- we identify the shallowest failing depth, generate a natural-language critique of the error, and use it to condition refined continuations that are grafted onto the tree as new branches. After re-scoring and reward propagation, previously uninformative trees become training-worthy, and grafted nodes participate in the standard \TreeRPO objective as regular children of their parent node.

Our contributions are:
\begin{itemize}
    \item We introduce the \emph{tree fitness function} $F(T)$, a zero-cost diagnostic that quantifies how much a tree will improve the LLM policy, and show that it correlates with per-tree policy gradient norms (H1: Spearman $r =$ \placeholder{X.XX}, $p <$ \placeholder{0.XX}).
    \item We propose \emph{critique-guided healing}, the first method to recover useful policy training signal from dead-wrong trees within the on-policy RL loop, and validate that it significantly increases tree fitness (H2: mean $\Delta F =$ \placeholder{X.XX}, $p <$ \placeholder{0.XX}).
    \item We integrate both mechanisms into \CATPO and demonstrate that the LLM policy trained with \CATPO improves over \TreeRPO and GRPO on four mathematical reasoning benchmarks, achieving \placeholder{XX.X}\% macro accuracy with Qwen2.5-Math-1.5B.
\end{itemize}


% ============================================================================
% 2. RELATED WORK
% ============================================================================
\section{Related Work}
\label{sec:related}

\subsection{Reinforcement Learning for LLM Reasoning}

Reinforcement learning from human feedback (RLHF)~\citep{rlhf} established the paradigm of optimizing LLM policy parameters with reward signals derived from human preferences. Proximal Policy Optimization (PPO)~\citep{ppo} has been the workhorse algorithm, though its reliance on a learned value function introduces instability and computational overhead.

Group Relative Policy Optimization (GRPO)~\citep{deepseekmathgrpo} addressed this by eliminating the value network entirely: for each prompt, a group of responses is sampled from the LLM policy, and the group-normalized reward serves as the advantage estimate for the policy gradient. GRPO and its variants form the backbone of recent reasoning breakthroughs including DeepSeek-R1~\citep{guo2025deepseek}. However, GRPO suffers from length bias and standard-deviation normalization bias~\citep{drgrpo}, and provides only trajectory-level advantage estimates. DAPO~\citep{dapo} introduced Dynamic Sampling (filtering zero-variance groups), token-level policy gradient loss, and decoupled clipping. VAPO~\citep{vapo} argued for value-based approaches with refined credit assignment for the policy. RLOO~\citep{rloo} and Open-Reasoner-Zero~\citep{openreasonerzero} demonstrated that simpler baselines can match or exceed GRPO under proper configuration.

\subsection{Tree-Based RL for Dense Reward Estimation}

The observation that tree-structured exploration improves reasoning has roots in inference-time methods: Tree-of-Thoughts~\citep{treeofthoughts}, RAP~\citep{rap}, and LATS~\citep{lats} use tree search to explore multiple reasoning paths without updating model weights. ReST-MCTS*~\citep{restmcts} and rStar-Math~\citep{rstarmath} extend this to training-time by collecting high-quality reasoning traces via MCTS for offline self-training or preference learning.

\TreeRPO~\citep{treerpo} was the first to integrate tree sampling directly into the on-policy RL loop for training the LLM policy: at each reasoning step, $N$ candidate continuations are sampled, forming an $N$-ary tree. Leaf outcomes are propagated bottom-up to estimate step-level advantages, and group-relative advantage computation is applied over sibling nodes. This provides dense, fine-grained advantage estimates for the policy gradient without training a separate reward model. TreeRL~\citep{treerl} further introduced entropy-guided branching (branching at high-uncertainty steps). Tree-GRPO~\citep{treegrpo} and Tree-OPO~\citep{treeopo} extended tree-based advantages to agent tasks and off-policy settings, respectively.

However, all existing tree-based RL methods treat every tree equally in the training objective. \CATPO is the first to \emph{diagnose} tree informativeness and \emph{intervene} on uninformative trees.

\subsection{Process Reward Models and Credit Assignment}

An alternative approach to step-level credit assignment trains a \emph{separate} neural network --- a Process Reward Model (PRM) --- to score intermediate reasoning steps. Let's Verify Step by Step~\citep{verify_stepbystep} established that process supervision (per-step correctness labels) outperforms outcome supervision for guiding the LLM policy, but requires expensive human annotation. Math-Shepherd~\citep{mathshepherd} and OmegaPRM~\citep{omegaprm} automated PRM data collection via Monte Carlo rollouts from intermediate states. PRIME~\citep{prime} computes token-level implicit process rewards through a log-ratio formulation without explicit PRM training. VinePPO~\citep{kazemnejad2025vineppo} replaces value networks with MC state-value estimates for policy optimization, and SPO~\citep{spo} provides segment-level credit assignment. These methods improve the quality of advantage estimates for the LLM policy gradient within existing rollout structures; \CATPO is complementary, operating at the tree level to filter and repair the rollouts themselves before they are used for policy updates.

\subsection{Sample Efficiency in LLM RL}

Wasted compute from uninformative rollouts is a recognized problem. DAPO's Dynamic Sampling~\citep{dapo} filters prompts where all responses succeed or fail. GRESO~\citep{greso} predicts uninformative prompts before rollout, achieving 2.4$\times$ speedup. CDAS~\citep{cdas} aligns problem difficulty to model competence. However, these methods operate at the \emph{prompt level}: they decide whether to train on a problem at all, but cannot diagnose the internal structure of a tree rollout. \CATPO operates at the \emph{tree level}, scoring and repairing the rollout structure itself.

\subsection{Self-Correction and Critique for LLMs}

Self-Refine~\citep{selfrefine} and Reflexion~\citep{reflexion} demonstrated that LLMs can iteratively improve outputs through self-critique at inference time. However, Huang et al.~\citep{huangselfcorrect} showed that intrinsic self-correction (without external feedback) often fails or degrades performance on reasoning tasks. SCoRe~\citep{score} addressed this by training self-correction ability through multi-turn online RL, proving that RL succeeds where prompting alone fails. GLoRe~\citep{glore} and RISE~\citep{rise} trained refinement models on synthetic correction trajectories. DRLC~\citep{drlc} used LLM-generated critiques as dense reward signals during RL training.

\CATPO differs from these approaches in a key structural way: rather than correcting complete trajectories or training a separate refinement model, we apply critique \emph{within} the tree structure at the specific point where reasoning fails, grafting corrected branches while preserving the shared prefix. This is more compute-efficient than regenerating entire solutions and naturally integrates with tree-based RL.


% ============================================================================
% 3. METHODOLOGY
% ============================================================================
\section{Methodology}
\label{sec:method}

We present \CATPO in four parts: background on tree-based RL (Section~\ref{sec:background}), the tree fitness function (Section~\ref{sec:fitness}), critique-guided healing (Section~\ref{sec:healing}), and the fitness-weighted training objective (Section~\ref{sec:objective}).

\subsection{Background: Tree-Based RL with GRPO}
\label{sec:background}

\paragraph{GRPO.} Given a prompt $q$, GRPO~\citep{deepseekmathgrpo} samples a group of $G$ complete responses $\{o_i\}_{i=1}^G$ from the current LLM policy $\pi_{\theta_\text{old}}$. Each response receives a reward $R_i$ from a rule-based verification function (e.g., checking whether the final answer matches the ground truth). The advantage for response $i$ is:
\begin{equation}
    \hat{A}_{i,t} = \frac{R_i - \text{mean}(\{R_j\}_{j=1}^G)}{\text{std}(\{R_j\}_{j=1}^G)}.
    \label{eq:grpo_adv}
\end{equation}

The LLM policy parameters $\theta$ are then updated via a clipped surrogate objective with KL regularization against a reference policy $\pi_\text{ref}$. Crucially, no separate value network or reward model is trained --- only the LLM policy is optimized.

\paragraph{TreeRPO.} Instead of sampling $G$ independent trajectories, \TreeRPO~\citep{treerpo} samples an $N$-ary tree of depth $D$ from the LLM policy. At each reasoning step, $N$ candidate continuations (each up to $L_\text{step}$ tokens) are generated, producing a tree where sibling nodes share a common prefix. Leaf nodes are scored by a verification function $\phi$, and rewards propagate bottom-up via mean aggregation:
\begin{equation}
    r_\text{node} = \mathbb{E}_{c \in \text{Children}(v)} [r_c].
    \label{eq:propagate}
\end{equation}

Advantages are computed over step-level sibling groups (children of the same parent), providing dense, fine-grained advantage estimates for the LLM policy gradient without training any auxiliary model.

\paragraph{Limitation.} Both GRPO and \TreeRPO treat all sampled data equally when updating the LLM policy. In GRPO, prompts where all responses succeed or fail produce zero-variance advantages. In \TreeRPO, this manifests at the tree level: dead-correct, dead-wrong, and stale trees contribute weak or zero gradients to the policy update. While DAPO's Dynamic Sampling~\citep{dapo} addresses this for GRPO by filtering zero-variance prompts, no analogous mechanism exists for tree-structured rollouts.


\subsection{Tree Fitness Function}
\label{sec:fitness}

We propose a \emph{tree fitness function} that quantifies how informative a tree $T$ is for improving the LLM policy. The function combines two signals available at zero extra cost after tree sampling:

\paragraph{Reward Entropy.} Let $\hat{p}$ be the fraction of leaves in $T$ that receive a positive reward:
\begin{equation}
    \mathcal{H}(\mathbf{r}) = -\big(\hat{p} \log_2 \hat{p} + (1 - \hat{p}) \log_2 (1 - \hat{p})\big).
    \label{eq:entropy}
\end{equation}
This binary entropy measures reward diversity. $\mathcal{H} = 0$ when all leaves agree (all correct or all incorrect), and $\mathcal{H} = 1$ when half succeed and half fail. A tree with zero entropy offers no contrastive signal for group-relative advantages.

\paragraph{Policy Surprise.} Let $\ell_n = \sum_t \log \pi_\theta(o_{n,t} \mid q, o_{n,<t})$ be the sum of token log-probabilities for non-root node $n$ under the current policy, and $r_n$ the propagated reward of node $n$. We compute:
\begin{equation}
    \rho_{\pi,r} = \text{Pearson}\big(\{\ell_n\}_{n \in T \setminus \text{root}},\ \{r_n\}_{n \in T \setminus \text{root}}\big).
    \label{eq:rho}
\end{equation}
When $\rho_{\pi,r} \approx 1$, the policy assigns high probability to high-reward paths and low probability to low-reward paths --- it has already internalized this tree's information. When $\rho_{\pi,r}$ is low or negative, the policy is \emph{surprised} by the reward distribution, indicating high learning potential.

\paragraph{Fitness Score.} The tree fitness is:
\begin{equation}
    \boxed{F(T) = \mathcal{H}(\mathbf{r}) \cdot (1 - \rho_{\pi,r})}
    \label{eq:fitness}
\end{equation}
with range $[0, 2]$. This multiplicative form ensures that both ingredients are necessary: a tree needs diverse outcomes \emph{and} policy surprise to be informative. All quantities are already computed during tree sampling and log-probability evaluation, so $F(T)$ incurs zero additional forward passes.

\paragraph{Tree Classification.} Based on $F(T)$, we classify trees into four regimes using thresholds $\tau_\text{low}$ and $\tau_\text{high}$:

\begin{table}[h]
\centering
\small
\begin{tabular}{llll}
\toprule
\textbf{Regime} & \textbf{Condition} & \textbf{Interpretation} & \textbf{Action} \\
\midrule
Dead-correct & $F \le \tau_\text{low}$, $\hat{p} > 0.5$ & All correct, nothing to learn & Downweight \\
Dead-wrong & $F \le \tau_\text{low}$, $\hat{p} \le 0.5$ & All wrong, no positive signal & Heal via critique \\
Stale & $\tau_\text{low} < F \le \tau_\text{high}$ & Policy already learned this & Moderate weight \\
Informative & $F > \tau_\text{high}$ & High learning potential & Full weight \\
\bottomrule
\end{tabular}
\caption{Tree classification regimes based on the fitness function $F(T)$.}
\label{tab:regimes}
\end{table}


\subsection{Critique-Guided Healing}
\label{sec:healing}

Dead-wrong trees ($F \approx 0$, all leaves incorrect) represent problems where the model consistently fails. Rather than discarding them, \CATPO attempts to \emph{heal} them by introducing corrected branches via natural-language critique.

\paragraph{Step 1: Identify the Failure Point.} We perform a breadth-first search from the root to find the shallowest depth $d^*$ where all children have propagated reward $r_c < \epsilon$ (with $\epsilon = 0.05$). This is the point where reasoning first goes entirely wrong.

\paragraph{Step 2: Generate a Critique.} The text from the root to the failing node at depth $d^*$ is extracted as a partial solution. A critic model (either the policy model itself via self-critique, or a separate model) is prompted to identify the specific error in the reasoning:

\vspace{0.5em}
\noindent\fbox{\parbox{0.95\linewidth}{
\small\textbf{Critique Prompt:} \textit{``You are given a math problem and a student's partial solution that leads to an incorrect answer. Identify the specific error in the reasoning and explain what went wrong.''}
}}
\vspace{0.5em}

The critique is generated at low temperature ($\tau = 0.3$) to produce a focused error diagnosis.

\paragraph{Step 3: Graft Refined Continuations.} The policy model is then prompted with the problem, the partial solution, and the generated critique, and asked to produce $k$ corrected continuations. These are added as new children of the node at depth $d^*$:

\vspace{0.5em}
\noindent\fbox{\parbox{0.95\linewidth}{
\small\textbf{Refinement Prompt:} \textit{``Given the following math problem, the student's incorrect attempt, and the critique identifying the error, provide a corrected solution continuing from the point of error.''}
}}
\vspace{0.5em}

\paragraph{Step 4: Re-score and Re-propagate.} The new leaves are scored by the verification function, rewards are re-propagated bottom-up via Eq.~\ref{eq:propagate}, and the tree fitness $F(T)$ is recomputed. If healing succeeds, some refined branches produce correct answers, converting a dead-wrong tree into one with positive reward diversity.

\paragraph{Hypothesis.} We hypothesize that critique-guided healing produces a statistically significant increase in tree fitness:
\begin{equation}
    \Delta F = F(T_\text{healed}) - F(T_\text{original}) > 0, \quad p < 0.05 \text{ (one-sided $t$-test)}.
\end{equation}


\subsection{CATPO Training Objective}
\label{sec:objective}

The full \CATPO algorithm proceeds in four phases per training iteration:

\paragraph{Phase 1: Tree Construction.} For each prompt $q$, sample an $N$-ary tree of depth $D$ following the \TreeRPO procedure. Score leaves with the verification function and propagate rewards bottom-up.

\paragraph{Phase 2: Fitness Diagnosis.} Compute $F(T)$ for each tree and classify into regimes (Table~\ref{tab:regimes}).

\paragraph{Phase 3: Critique Healing.} For each dead-wrong tree: identify the shallowest failure point, generate a critique, graft $k$ refined continuations, re-score, and re-propagate rewards.

\paragraph{Phase 4: Fitness-Weighted Policy Update.} After healing, all nodes --- both original and grafted --- are treated uniformly. We weight each tree's contribution to the policy gradient by its normalized fitness:
\begin{equation}
    w(T) = \frac{F(\tilde{T})}{\sum_{T' \in \mathcal{B}} F(\tilde{T'})},
    \label{eq:fitness_weight}
\end{equation}
where $\tilde{T}$ denotes the tree after healing (if applicable). This concentrates the LLM policy gradient on informative trees and downweights dead-correct and stale ones.

\paragraph{Objective.} We adopt the same clipped surrogate objective as \TreeRPO, with fitness weighting applied per tree:
\begin{equation}
\begin{aligned}
    \mathcal{J}_\text{CATPO}(\theta) = \sum_{T \in \mathcal{B}} w(T) \sum_{n \in T} \frac{1}{|o_n|} \sum_{t=1}^{|o_n|} \Big(
    &\min\big(r_{n,t}(\theta)\, \hat{A}_{n,t},\ \text{clip}(r_{n,t}(\theta), 1-\varepsilon, 1+\varepsilon)\, \hat{A}_{n,t}\big) \\
    &- \beta\, D_\text{KL}(\pi_\theta \| \pi_\text{ref}) \Big),
\end{aligned}
\label{eq:catpo_loss}
\end{equation}
where $r_{n,t}(\theta) = \pi_\theta(o_{n,t} \mid q, o_{n,<t}) / \pi_{\theta_\text{old}}(o_{n,t} \mid q, o_{n,<t})$ is the standard importance ratio and $\hat{A}_{n,t}$ is the group-relative advantage computed over sibling nodes following \TreeRPO. The summation over $n \in T$ includes both original and grafted nodes; the only difference from \TreeRPO is the per-tree fitness weight $w(T)$.

The key insight is that critique healing does not require a modified loss function. Grafted nodes become regular children of their parent node after re-propagation, so the standard \TreeRPO advantage computation and clipped objective apply directly. The fitness weight ensures that successfully healed trees (which now have diverse leaf outcomes and higher $F$) receive appropriately larger gradient contributions.

The complete algorithm is summarized in Algorithm~\ref{alg:catpo}.

\begin{algorithm}[t]
\caption{\CATPO: Critique-Augmented Tree Policy Optimization}
\label{alg:catpo}
\begin{algorithmic}[1]
\REQUIRE LLM policy $\pi_\theta$, reference policy $\pi_\text{ref}$, training data $\mathcal{D}$, branching factor $N$, max depth $D$, refinements $k$, thresholds $\tau_\text{low}, \tau_\text{high}$
\FOR{each training iteration}
    \STATE Sample batch of prompts $\{q_i\}$ from $\mathcal{D}$
    \STATE \textbf{Phase 1:} For each $q_i$, sample $N$-ary tree $T_i$ of depth $D$; score leaves; propagate rewards
    \STATE \textbf{Phase 2:} Compute $F(T_i) = \mathcal{H}(\mathbf{r}_i) \cdot (1 - \rho_{\pi,r_i})$; classify each tree
    \STATE \textbf{Phase 3:} For each dead-wrong tree $T_i$ (where $F(T_i) \le \tau_\text{low}$ and $\hat{p}_i \le 0.5$):
    \STATE \quad Find shallowest failing depth $d^*$
    \STATE \quad Generate critique of partial solution up to $d^*$
    \STATE \quad Generate $k$ refined continuations conditioned on critique
    \STATE \quad Graft onto tree; re-score leaves; re-propagate rewards
    \STATE \textbf{Phase 4:} Compute fitness weights $w(T_i) = F(\tilde{T}_i) / \sum_j F(\tilde{T}_j)$
    \STATE Update $\theta$ via $\nabla_\theta \mathcal{J}_\text{CATPO}(\theta)$ \hfill (Eq.~\ref{eq:catpo_loss})
\ENDFOR
\end{algorithmic}
\end{algorithm}


% ============================================================================
% 4. EXPERIMENTS
% ============================================================================
\section{Experiments}
\label{sec:experiments}

\subsection{Experimental Setup}

\paragraph{Datasets.} We follow the \TreeRPO experimental protocol. For training, we use the training split of the MATH dataset~\citep{hendrycksmath2021}, which contains 7.5k high-quality mathematical reasoning problems. For evaluation, we report results on four widely used benchmarks: MATH-500~\citep{verify_stepbystep} (500 representative problems from the MATH test split), OlympiadBench~\citep{olympiadbench}, MinervaMath~\citep{minerva}, and AIME24.

\paragraph{Base Models.} Our experiments use the Qwen2.5-Math series~\citep{qwen25math}. We primarily report results on Qwen2.5-Math-1.5B, with scaling experiments on Qwen2.5-Math-7B.

\paragraph{Baselines.} We compare against:
\begin{itemize}
    \item \textbf{GRPO}~\citep{deepseekmathgrpo}: flat trajectory sampling with group-relative advantages.
    \item \textbf{\TreeRPO}~\citep{treerpo}: tree-structured sampling with step-level group-relative advantages (our direct base method).
\end{itemize}

\paragraph{Implementation Details.} We build on the rllm~\citep{deepcoder2025,deepscaler2025} framework derived from verl~\citep{verl}, with vLLM~\citep{vllm} for efficient inference. For GRPO: temperature 0.6, 8 rollouts per prompt, batch size 128, learning rate 1e-6, max response length 1152. For \TreeRPO and \CATPO: temperature 0.6, branching factor $N=8$, max depth $D=3$, max step length $L_\text{step}=384$, pruning threshold $\tau=0.1$. \CATPO-specific: $\tau_\text{low}=0.2$, $\tau_\text{high}=0.8$, $k=4$ refinements. KL coefficient $\beta=0.001$, entropy loss coefficient $\alpha=-0.001$. All experiments are conducted on \placeholder{N}$\times$A800 GPUs.

\paragraph{Metrics.} We report \textbf{Pass@1 (Avg@8)} accuracy: for each test problem, 8 responses are sampled at temperature 0.6, and accuracy is averaged across samples. We report individual benchmark scores and the \textbf{Macro Accuracy} (unweighted average across four benchmarks).


\subsection{Hypothesis Validation}

Before evaluating \CATPO end-to-end, we validate the two core hypotheses underlying our approach.

\paragraph{H1: Fitness Predicts Gradient Utility.}

We build trees offline for \placeholder{N} problems from the MATH training set using Qwen2.5-Math-1.5B, compute $F(T)$ for each tree, and measure the per-tree gradient norm $\|g(T)\|^2$ via REINFORCE-style forward-backward passes. If $F(T)$ is a good proxy for training utility, it should correlate positively with $\|g(T)\|^2$.

\begin{table}[h]
\centering
\begin{tabular}{lcc}
\toprule
\textbf{Correlation Metric} & \textbf{Value} & \textbf{$p$-value} \\
\midrule
Pearson $r$ & \placeholder{X.XX} & \placeholder{X.XX} \\
Spearman $\rho$ & \placeholder{X.XX} & \placeholder{X.XX} \\
\bottomrule
\end{tabular}
\caption{Correlation between tree fitness $F(T)$ and per-tree gradient norm $\|g(T)\|^2$ on \placeholder{N} trees from MATH training set. Target: Spearman $\rho > 0.4$, $p < 0.01$.}
\label{tab:h1}
\end{table}

\placeholder{Discussion of H1 results: do we meet the target? What does the scatter plot look like? Are there outliers?}

\paragraph{H2: Critique Heals Dead-Wrong Trees.}

From the same tree set, we select all dead-wrong trees ($F \le 0.2$, $\hat{p} \le 0.5$) and apply critique-guided healing with $k=4$ refinements. We measure the fitness improvement $\Delta F = F_\text{after} - F_\text{before}$.

\begin{table}[h]
\centering
\begin{tabular}{lccc}
\toprule
& \textbf{$F_\text{before}$} & \textbf{$F_\text{after}$} & \textbf{$\Delta F$} \\
\midrule
Mean & \placeholder{X.XX} & \placeholder{X.XX} & \placeholder{X.XX} \\
Std & \placeholder{X.XX} & \placeholder{X.XX} & \placeholder{X.XX} \\
\midrule
One-sided $t$-test $p$-value & \multicolumn{3}{c}{\placeholder{X.XX}} \\
\% trees with $\Delta F > 0$ & \multicolumn{3}{c}{\placeholder{XX}\%} \\
\bottomrule
\end{tabular}
\caption{Critique healing results on \placeholder{N} dead-wrong trees. Target: mean $\Delta F > 0$, $p < 0.05$.}
\label{tab:h2}
\end{table}

\placeholder{Discussion of H2 results: healing success rate, typical failure modes, effect of self-critique vs. separate critic model.}


\subsection{Main Results}

\begin{table}[t]
    \centering
    \renewcommand\arraystretch{1.3}
    \setlength{\tabcolsep}{5pt}
    \caption{Overall \textit{Pass@1} (\textit{Avg@8}) performance on four mathematical reasoning benchmarks.}
    \begin{tabular}{l|c|c|c|c|c}
       \toprule
       \textbf{Method} & \textbf{AIME24} & \textbf{MATH-500} & \textbf{OlympiadBench} & \textbf{Minerva} & \textbf{Macro Acc.} \\
       \midrule
       \rowcolor{gray!16} \multicolumn{6}{c}{\textit{Qwen2.5-Math-1.5B as the Base Model}} \\
        GRPO & 13.8 & 67.9 & 28.5 & 20.5 & 32.7 \\
        \TreeRPO & 16.8 {\scriptsize\color{darkgreen}$\uparrow$3.0} & 70.7 {\scriptsize\color{darkgreen}$\uparrow$2.8} & 30.9 {\scriptsize\color{darkgreen}$\uparrow$2.4} & 24.0 {\scriptsize\color{darkgreen}$\uparrow$3.5} & 35.6 {\scriptsize\color{darkgreen}$\uparrow$2.9} \\
        \textbf{\CATPO} & \placeholder{XX.X} {\scriptsize\color{darkgreen}$\uparrow$\placeholder{X.X}} & \placeholder{XX.X} {\scriptsize\color{darkgreen}$\uparrow$\placeholder{X.X}} & \placeholder{XX.X} {\scriptsize\color{darkgreen}$\uparrow$\placeholder{X.X}} & \placeholder{XX.X} {\scriptsize\color{darkgreen}$\uparrow$\placeholder{X.X}} & \placeholder{XX.X} {\scriptsize\color{darkgreen}$\uparrow$\placeholder{X.X}} \\
        \midrule
        \rowcolor{gray!16} \multicolumn{6}{c}{\textit{Qwen2.5-Math-7B as the Base Model}} \\
        GRPO & 26.7 & 74.3 & 34.7 & 27.1 & 40.7 \\
        \TreeRPO & 26.7 & 75.5 {\scriptsize\color{darkgreen}$\uparrow$1.2} & 35.4 {\scriptsize\color{darkgreen}$\uparrow$0.7} & 28.1 {\scriptsize\color{darkgreen}$\uparrow$1.0} & 41.4 {\scriptsize\color{darkgreen}$\uparrow$0.7} \\
        \textbf{\CATPO} & \placeholder{XX.X} & \placeholder{XX.X} & \placeholder{XX.X} & \placeholder{XX.X} & \placeholder{XX.X} {\scriptsize\color{darkgreen}$\uparrow$\placeholder{X.X}} \\
        \bottomrule
    \end{tabular}
    \label{tab:main}
\end{table}

Table~\ref{tab:main} presents the main results. For Qwen2.5-Math-1.5B, \CATPO achieves a Macro Accuracy of \placeholder{XX.X}\%, improving over \TreeRPO by \placeholder{+X.X}\% and over GRPO by \placeholder{+X.X}\%. \placeholder{Discuss per-benchmark trends: where does CATPO help most? Is the improvement larger on harder benchmarks (AIME, OlympiadBench) where dead-wrong trees are more common?}

For Qwen2.5-Math-7B, \placeholder{discuss whether the improvement persists at larger scale, and whether the MATH training data difficulty limitation noted by TreeRPO authors affects CATPO similarly.}


\subsection{Training Dynamics}

\placeholder{Include training curves (accuracy vs. training steps) for GRPO, TreeRPO, and CATPO on each benchmark. Describe:
\begin{itemize}
    \item Whether CATPO converges faster (sample efficiency).
    \item Whether the improvement is consistent across training or mainly in early/late phases.
    \item Training stability (any signs of entropy collapse or reward hacking).
\end{itemize}
Reference figures similar to TreeRPO's Figure 3 (performance-bs128.pdf).}


\subsection{Token Efficiency}

\begin{table}[h]
    \centering
    \renewcommand\arraystretch{1.2}
    \caption{Average response length (tokens) on test benchmarks. Lower is better at equal accuracy.}
    \begin{tabular}{l|c|c|c|c|c}
    \toprule
    \textbf{Method} & \textbf{AIME24} & \textbf{MATH-500} & \textbf{OlympiadBench} & \textbf{Minerva} & \textbf{Avg. Length} \\
    \midrule
    GRPO & \placeholder{XXX} & \placeholder{XXX} & \placeholder{XXX} & \placeholder{XXX} & \placeholder{XXX} \\
    \TreeRPO & \placeholder{XXX} & \placeholder{XXX} & \placeholder{XXX} & \placeholder{XXX} & \placeholder{XXX} \\
    \textbf{\CATPO} & \placeholder{XXX} & \placeholder{XXX} & \placeholder{XXX} & \placeholder{XXX} & \placeholder{XXX} \\
    \bottomrule
    \end{tabular}
    \label{tab:length}
\end{table}

\TreeRPO demonstrated an 18.1\% reduction in average response length compared to GRPO. \placeholder{Report CATPO's response length and discuss whether fitness weighting (which downweights stale trees where the model tends to generate verbose but already-known reasoning) further reduces token usage.}


\subsection{Ablation Study}

To disentangle the contributions of fitness weighting and critique healing, we evaluate three ablation variants:

\begin{table}[h]
    \centering
    \renewcommand\arraystretch{1.2}
    \caption{Ablation study on Qwen2.5-Math-1.5B. Macro Accuracy on four benchmarks.}
    \begin{tabular}{l|cc|c}
    \toprule
    \textbf{Method} & \textbf{Fitness Weighting} & \textbf{Critique Healing} & \textbf{Macro Acc.} \\
    \midrule
    \TreeRPO (baseline) & \texttimes & \texttimes & 35.6 \\
    + Fitness weighting only & \checkmark & \texttimes & \placeholder{XX.X} \\
    + Critique healing only & \texttimes & \checkmark & \placeholder{XX.X} \\
    \textbf{\CATPO (full)} & \checkmark & \checkmark & \placeholder{XX.X} \\
    \bottomrule
    \end{tabular}
    \label{tab:ablation}
\end{table}

\placeholder{Discuss ablation results:
\begin{itemize}
    \item Does fitness weighting alone improve over TreeRPO? (Expected: yes, by focusing gradients on informative trees.)
    \item Does critique healing alone improve? (Expected: yes, by recovering signal from dead-wrong trees.)
    \item Is the full combination better than either alone? (Expected: yes, since they target different failure modes.)
\end{itemize}}


\subsection{Analysis of Tree Fitness Distribution}

\begin{table}[h]
    \centering
    \renewcommand\arraystretch{1.2}
    \caption{Distribution of tree regimes during training (averaged over all iterations).}
    \begin{tabular}{l|cccc}
    \toprule
    \textbf{Method} & \textbf{Dead-Correct} & \textbf{Dead-Wrong} & \textbf{Stale} & \textbf{Informative} \\
    \midrule
    \TreeRPO & \placeholder{XX}\% & \placeholder{XX}\% & \placeholder{XX}\% & \placeholder{XX}\% \\
    \CATPO (before healing) & \placeholder{XX}\% & \placeholder{XX}\% & \placeholder{XX}\% & \placeholder{XX}\% \\
    \CATPO (after healing) & \placeholder{XX}\% & \placeholder{XX}\% & \placeholder{XX}\% & \placeholder{XX}\% \\
    \bottomrule
    \end{tabular}
    \label{tab:regimes_dist}
\end{table}

\placeholder{Discuss: What fraction of trees are dead-wrong and thus eligible for healing? How many are successfully healed (converted from dead-wrong to informative or stale)? Does the distribution shift during training as the model improves?}


\subsection{Hyperparameter Sensitivity}

\paragraph{Batch Size.} Following \TreeRPO, we evaluate \CATPO with batch sizes 16 and 128.

\begin{table}[h]
    \centering
    \begin{tabular}{l|cc}
    \toprule
    \textbf{Method} & \textbf{bsz=16} & \textbf{bsz=128} \\
    \midrule
    GRPO & \placeholder{XX.X} & 32.7 \\
    \TreeRPO & \placeholder{XX.X} & 35.6 \\
    \textbf{\CATPO} & \placeholder{XX.X} & \placeholder{XX.X} \\
    \bottomrule
    \end{tabular}
    \caption{Macro Accuracy under different batch sizes (Qwen2.5-Math-1.5B).}
    \label{tab:batchsize}
\end{table}

\placeholder{Discuss robustness to batch size.}

\paragraph{Number of Refinements $k$.} We vary the number of critique-conditioned refinements grafted onto dead-wrong trees.

\begin{table}[h]
    \centering
    \begin{tabular}{l|cccc}
    \toprule
    $k$ & 1 & 2 & 4 & 8 \\
    \midrule
    Macro Acc. & \placeholder{XX.X} & \placeholder{XX.X} & \placeholder{XX.X} & \placeholder{XX.X} \\
    Healing Success Rate & \placeholder{XX}\% & \placeholder{XX}\% & \placeholder{XX}\% & \placeholder{XX}\% \\
    \bottomrule
    \end{tabular}
    \caption{Effect of the number of refinements $k$ on \CATPO performance.}
    \label{tab:k_ablation}
\end{table}

\placeholder{Discuss: is there diminishing returns? What is the compute-performance tradeoff?}

\paragraph{Fitness Thresholds $\tau_\text{low}, \tau_\text{high}$.} \placeholder{Report sensitivity to the classification thresholds. Are the results robust to reasonable variations?}


% ============================================================================
% 5. CONCLUSION
% ============================================================================
\section{Conclusion}
\label{sec:conclusion}

We introduced \CATPO, a method that enhances tree-based reinforcement learning for training LLM reasoning policies through two complementary mechanisms: a zero-cost tree fitness function that identifies which trees are most informative for the policy gradient, and a critique-guided healing procedure that recovers useful training signal from dead-wrong trees. Our fitness function $F(T) = \mathcal{H}(\mathbf{r}) \cdot (1 - \rho_{\pi,r})$ combines reward diversity with policy surprise to quantify each tree's potential for improving the LLM, and we validated that it correlates with per-tree policy gradient norms (H1) and that critique healing significantly increases the fitness of failing trees (H2).

By weighting the policy gradient loss by tree fitness, \CATPO concentrates the LLM's parameter updates on the most informative signal while rescuing otherwise wasted compute through critique healing. No additional reward model or value network is trained --- \CATPO, like \TreeRPO, optimizes only the LLM policy parameters. Experiments on Qwen2.5-Math-1.5B demonstrate that the LLM policy trained with \CATPO improves over both GRPO (\placeholder{+X.X}\%) and \TreeRPO (\placeholder{+X.X}\%) in macro accuracy across four mathematical reasoning benchmarks.

\paragraph{Limitations and Future Work.}
\begin{itemize}
    \item \textbf{Critique quality.} When the policy model serves as its own critic, self-critique quality is bounded by the model's own capabilities. Scaling to separate, stronger critic models may yield further gains but introduces additional inference cost.
    \item \textbf{Correlation noise.} The Pearson correlation $\rho_{\pi,r}$ can be noisy with small trees (few nodes). Exploring rank-based or regularized alternatives may improve robustness.
    \item \textbf{Scale.} Our experiments are primarily on 1.5B parameters due to compute constraints. We plan to evaluate on 7B and 32B models, where the fraction of informative trees may differ.
    \item \textbf{Beyond math.} While we focus on mathematical reasoning with rule-based verification, \CATPO's fitness function generalizes to any domain with verifiable rewards (e.g., code generation). Extending to code benchmarks is a natural next step.
    \item \textbf{Engineering optimization.} Tree sampling with critique healing adds overhead relative to \TreeRPO. KV cache sharing across the tree's shared prefixes~\citep{vllm} and batched critique generation could significantly reduce wall-clock time.
\end{itemize}


% ============================================================================
% ACKNOWLEDGMENTS
% ============================================================================
\section*{Acknowledgments}
\placeholder{Acknowledgments.}


% ============================================================================
% BIBLIOGRAPHY
% ============================================================================
\bibliography{colm2026_conference}
\bibliographystyle{colm2026_conference}


% ============================================================================
% APPENDIX
% ============================================================================
\appendix

\section{Detailed Fitness Function Derivation}
\label{app:fitness_derivation}

\paragraph{Motivation.} We seek a scalar function $F: \mathcal{T} \to \mathbb{R}_{\ge 0}$ that ranks trees by their expected contribution to policy improvement. A useful tree should satisfy two properties: (1) its leaf outcomes are diverse (providing contrastive signal for group-relative advantages), and (2) the policy does not already predict which branches succeed (ensuring the tree carries novel information).

\paragraph{Entropy Term.} The first property is captured by the binary entropy of the leaf success rate $\hat{p} = \frac{1}{|\mathcal{L}|} \sum_{\ell \in \mathcal{L}} \mathbf{1}[r_\ell > 0]$:
\begin{equation}
    \mathcal{H}(\mathbf{r}) = -\hat{p} \log_2 \hat{p} - (1 - \hat{p}) \log_2(1 - \hat{p}).
\end{equation}
$\mathcal{H} \in [0, 1]$, with $\mathcal{H} = 0$ when all leaves agree and $\mathcal{H} = 1$ when outcomes are maximally diverse.

\paragraph{Surprise Term.} The second property is captured by the Pearson correlation between per-node log-probabilities and propagated rewards:
\begin{equation}
    \rho_{\pi,r} = \frac{\sum_{n} (\ell_n - \bar{\ell})(r_n - \bar{r})}{\sqrt{\sum_n (\ell_n - \bar{\ell})^2} \sqrt{\sum_n (r_n - \bar{r})^2}},
\end{equation}
where $\ell_n = \sum_t \log \pi_\theta(o_{n,t} | q, o_{n,<t})$ and $r_n$ is the propagated reward. $\rho \in [-1, 1]$, with $\rho \approx 1$ indicating the policy already knows the reward structure.

\paragraph{Combined Score.} The multiplicative form $F(T) = \mathcal{H}(\mathbf{r}) \cdot (1 - \rho_{\pi,r})$ ensures both properties are necessary. The range is $[0, 2]$: $F = 0$ when either $\mathcal{H} = 0$ (no diversity) or $\rho = 1$ (no surprise), and $F$ can reach 2 when $\mathcal{H} = 1$ and $\rho = -1$ (maximal diversity and anti-correlated policy confidence).

\paragraph{Computational Cost.} All ingredients --- leaf rewards, propagated rewards, and token log-probabilities --- are already computed during tree sampling and the log-probability forward pass. Computing $F(T)$ adds only $O(|T|)$ arithmetic operations per tree, which is negligible compared to the generation cost.


\section{Critique Prompt Templates}
\label{app:prompts}

\paragraph{Critique Generation.}
\begin{verbatim}
System: You are a mathematical reasoning critic. Your job is to
identify specific errors in student solutions.

User: Problem: {problem_text}

The student wrote the following partial solution, but it leads
to an incorrect final answer:

{partial_solution}

Identify the specific mathematical or logical error in the
reasoning above. Be precise about which step is wrong and why.
\end{verbatim}

\paragraph{Refined Continuation.}
\begin{verbatim}
System: You are a careful mathematical problem solver. You are
given a problem, a student's incorrect attempt, and a critique
identifying the error. Provide a corrected continuation.

User: Problem: {problem_text}

Incorrect attempt: {partial_solution}

Critique: {critique_text}

Now provide a corrected solution, continuing from the point
where the error was identified. Show your work step by step.
\end{verbatim}


\section{Extended Results}
\label{app:extended}

\subsection{Per-Benchmark Training Curves}

\placeholder{Include per-benchmark training curves (Pass@1 vs. training steps) for GRPO, TreeRPO, and CATPO on MATH-500, OlympiadBench, MinervaMath, and AIME24. Analogous to TreeRPO's Figures 3 and 4.}

\subsection{Response Length Analysis}

\placeholder{Include per-benchmark response length curves (tokens vs. training steps) for all methods. Analogous to TreeRPO's Figure 5.}

\subsection{Fitness Distribution Over Training}

\placeholder{Show how the distribution of tree regimes (dead-correct, dead-wrong, stale, informative) evolves over training steps for both TreeRPO and CATPO. This should show that as the model improves, more trees become dead-correct (model solves them) and fewer are informative (the frontier shrinks), motivating the need for fitness-weighted training.}

\subsection{Critique Quality Examples}

\placeholder{Include 2-3 examples showing: (1) a dead-wrong tree's failing partial solution, (2) the generated critique, (3) the refined continuation, and (4) whether it leads to a correct answer. Include both a success case and a failure case.}


\section{Computational Overhead Analysis}
\label{app:compute}

\placeholder{Report wall-clock time per training step for GRPO, TreeRPO, and CATPO. Break down CATPO's overhead into: (1) fitness computation (expected: negligible), (2) critique generation (the main cost), and (3) refinement generation + re-scoring. Report the fraction of trees that require healing per iteration and the amortized cost.}


\end{document}
